\chapter{Linear Algebra}

\section{Complex Numbers and Fields}

\begin{definition}[Complex numbers]
A complex number is an ordered pair $(a,b)$ of reals, written $a+bi$; $\mathbb{C}$ is the set of all of them, with
\[
    (a+bi)+(c+di) = (a+c)+(b+d)i,\qquad (a+bi)(c+di) = (ac-bd)+(ad+bc)i .
\]
\end{definition}
Complex arithmetic is commutative, associative and distributive; it has identities ($\lambda+0=\lambda$, $\lambda\cdot1=\lambda$), a unique additive inverse $-\alpha$ for every $\alpha$, and a unique multiplicative inverse $1/\alpha$ for every $\alpha\neq0$. Subtraction and division are defined through these: $\beta-\alpha = \beta+(-\alpha)$ and $\beta/\alpha = \beta\cdot(1/\alpha)$. Any set with these properties is a \emph{field}; below $\mathbb{F}$ denotes $\mathbb{R}$ or $\mathbb{C}$, and its elements are called \emph{scalars}.

\begin{definition}[Lists and $\mathbb{F}^n$]
A list of length $n$ is an ordered collection of $n$ elements; two lists are equal iff they have the same length and the same elements in the same order. $\mathbb{F}^n = \{(x_1,\ldots,x_n) : x_k\in\mathbb{F}\}$, with $x_k$ the $k$th coordinate. Addition and scalar multiplication are componentwise,
\[
    (x_1,\ldots,x_n)+(y_1,\ldots,y_n) = (x_1+y_1,\ldots,x_n+y_n),\qquad \lambda(x_1,\ldots,x_n) = (\lambda x_1,\ldots,\lambda x_n),
\]
with zero $0 = (0,\ldots,0)$ and additive inverse $-x = (-x_1,\ldots,-x_n)$.
\end{definition}

\section{Vector Spaces}

\begin{definition}[Vector space]
A vector space over $\mathbb{F}$ is a set $V$ with an addition $V\times V\to V$ and a scalar multiplication $\mathbb{F}\times V\to V$ such that, for all $x,y,z\in V$ and $a,b\in\mathbb{F}$:
\begin{enumerate}
    \item commutativity: $x+y = y+x$;
    \item associativity: $(x+y)+z = x+(y+z)$ and $(ab)x = a(bx)$;
    \item additive identity: there is $0\in V$ with $x+0 = x$;
    \item additive inverse: for each $x$ there is $w$ with $x+w = 0$;
    \item multiplicative identity: $1x = x$;
    \item distributivity: $a(x+y) = ax+ay$ and $(a+b)x = ax+bx$.
\end{enumerate}
One must say ``over $\mathbb{F}$'', since scalar multiplication depends on the field.
\end{definition}

\begin{example}[Function spaces]
For any set $S$, the set $\mathbb{F}^S$ of functions $S\to\mathbb{F}$ is a vector space under $(f+g)(x) = f(x)+g(x)$ and $(\lambda f)(x) = \lambda f(x)$. $\mathbb{F}^n$ is the special case $S = \{1,\ldots,n\}$, with $x_k = x(k)$.
\end{example}

\begin{theorem}[Elementary properties]
In any vector space:
\begin{enumerate}
    \item the additive identity is unique;
    \item each element has a unique additive inverse, written $-v$ (and $v-w \equiv v+(-w)$);
    \item $0v = 0$ for every $v\in V$, and $a0 = 0$ for every $a\in\mathbb{F}$;
    \item $(-1)v = -v$.
\end{enumerate}
\end{theorem}
\begin{proof}
(1) If $0,0'$ are both identities, $0' = 0'+0 = 0+0' = 0$. (2) If $w,w'$ are inverses of $v$, then $w = w+(v+w') = (w+v)+w' = w'$. (3) $0v = (0+0)v = 0v+0v$; adding $-0v$ to both sides gives $0v = 0$, and similarly $a0 = a(0+0) = a0+a0$. (4) $v+(-1)v = (1-1)v = 0v = 0$.
\end{proof}

\begin{example}[Non-examples]
$\mathbb{R}\cup\{\pm\infty\}$ fails associativity: $(-\infty+\infty)+\infty = \infty$ but $-\infty+(\infty+\infty) = 0$. The integer lattice $\mathbb{Z}^2$ is not a vector space over $\mathbb{R}$, since $\lambda x\notin\mathbb{Z}^2$ for $\lambda\notin\mathbb{Z}$. Neither is $\mathbb{R}^2$ viewed as a subset of $\mathbb{C}^2$ over $\mathbb{C}$, since multiplying by $i$ leaves it.
\end{example}

\section{Subspaces, Sums and Direct Sums}

\begin{theorem}[Conditions for a subspace]
A subset $U\subseteq V$ is a subspace (a vector space with the same operations) iff
\begin{enumerate}
    \item $0\in U$;
    \item $u,w\in U\Rightarrow u+w\in U$;
    \item $a\in\mathbb{F}$, $u\in U\Rightarrow au\in U$.
\end{enumerate}
\end{theorem}
The remaining axioms are inherited from $V$, and closure under scalars gives $-u = (-1)u\in U$. For instance, the subspaces of $\mathbb{R}^2$ are $\{0\}$, all lines through the origin, and $\mathbb{R}^2$. The functions on $[0,1]$ with $\int_0^1 f = b$ form a subspace iff $b = 0$; so do the differentiable functions on $[-4,4]$ with $f'(-1) = 3f(2)$.

\begin{definition}[Sum of subspaces]
$V_1+\cdots+V_m = \{v_1+\cdots+v_m : v_k\in V_k\}$. It is the smallest subspace of $V$ containing $V_1,\ldots,V_m$, the analogue of the union of sets.
\end{definition}

\begin{definition}[Direct sum]
$V_1+\cdots+V_m$ is a \emph{direct sum}, written $V_1\oplus\cdots\oplus V_m$, if every element of it can be written in only one way as $v_1+\cdots+v_m$ with $v_k\in V_k$.
\end{definition}

\begin{theorem}[Condition for a direct sum]
$V_1+\cdots+V_m$ is direct iff the only way to write $0 = v_1+\cdots+v_m$ with $v_k\in V_k$ is $v_k = 0$ for all $k$. Two subspaces give a direct sum, $U+W = U\oplus W$, iff $U\cap W = \{0\}$.
\end{theorem}
\begin{proof}
If the sum is direct, $0 = 0+\cdots+0$ is the only representation. Conversely, if $x = \sum v_k = \sum v'_k$, subtracting gives $0 = \sum(v_k - v'_k)$, so $v_k = v'_k$ for all $k$.
\end{proof}

\begin{example}[A sum that is not direct]
In $\mathbb{F}^3$, let $V_1 = \{(x,y,0)\}$, $V_2 = \{(0,0,z)\}$, $V_3 = \{(0,y,y)\}$. Then $\mathbb{F}^3 = V_1+V_2+V_3$, but the sum is not direct: $(0,0,0) = (0,1,0)+(0,0,1)+(0,-1,-1)$.
\end{example}

\begin{fact}[Unions and intersections]
The intersection of subspaces is a subspace. A union $V_1\cup V_2$ is a subspace iff one contains the other: if neither does, take $v_1\in V_1\setminus V_2$ and $v_2\in V_2\setminus V_1$; then $v_1+v_2$ lies in neither. Also $U+U = U$.
\end{fact}

\section{Span, Linear Independence and Bases}

\begin{definition}[Span]
A \emph{linear combination} of $v_1,\ldots,v_m$ is $a_1v_1+\cdots+a_mv_m$ with $a_i\in\mathbb{F}$; $\operatorname{span}(v_1,\ldots,v_m)$ is the set of all of them. It is the smallest subspace containing $v_1,\ldots,v_m$, since any subspace containing the list contains all its linear combinations. If the span equals $V$, the list \emph{spans} $V$; $V$ is \emph{finite-dimensional} if some finite list spans it.
\end{definition}

\begin{definition}[Polynomials]
$p:\mathbb{F}\to\mathbb{F}$ is a polynomial of degree $m$ if $p(z) = a_0 + a_1z+\cdots+a_mz^m$ with $a_m\neq0$ (the zero polynomial has degree $-\infty$). $\mathcal{P}(\mathbb{F})$ is a subspace of $\mathbb{F}^{\mathbb{F}}$, and $\mathcal{P}_m(\mathbb{F}) = \operatorname{span}(1,z,\ldots,z^m)$, the polynomials of degree at most $m$, is finite-dimensional. The coefficients of a polynomial are uniquely determined by the function.
\end{definition}

\begin{definition}[Linear independence]
$v_1,\ldots,v_m$ are \emph{linearly independent} if $a_1v_1+\cdots+a_mv_m = 0$ only for $a_1 = \cdots = a_m = 0$; equivalently, every vector in their span has exactly one representation as a linear combination of them. Otherwise they are \emph{linearly dependent}.
\end{definition}

\begin{theorem}[Linear dependence lemma]
If $v_1,\ldots,v_m$ are linearly dependent, there is a $k$ with $v_k\in\operatorname{span}(v_1,\ldots,v_{k-1})$, and removing $v_k$ leaves the span unchanged.
\end{theorem}
\begin{proof}
Take $a_1v_1+\cdots+a_mv_m = 0$ with not all $a_i = 0$, and let $k$ be the largest index with $a_k\neq0$. Then $v_k = -\frac{a_1}{a_k}v_1-\cdots-\frac{a_{k-1}}{a_k}v_{k-1}$. Any $v = \sum b_iv_i$ in the original span can then be rewritten without $v_k$, so the span is unchanged. (For $k=1$, this says $v_1 = 0$.)
\end{proof}

\begin{theorem}
In a finite-dimensional vector space, the length of every linearly independent list is at most the length of every spanning list. Every subspace of a finite-dimensional space is finite-dimensional.
\end{theorem}

\begin{definition}[Basis]
A basis of $V$ is a linearly independent list that spans $V$. Equivalently, $v_1,\ldots,v_n$ is a basis iff every $v\in V$ can be written uniquely as $v = a_1v_1+\cdots+a_nv_n$.
\end{definition}

\begin{theorem}
If $V$ is finite-dimensional and $U$ is a subspace, there is a subspace $W$ with $V = U\oplus W$.
\end{theorem}

\section{Linear Maps and Matrices}

\begin{definition}[Linear map]
$T:V\to W$ is linear if $T(u+v) = Tu+Tv$ and $T(\lambda v) = \lambda Tv$ for all $u,v\in V$ and $\lambda\in\mathbb{F}$. $\mathcal{L}(V,W)$ denotes the set of linear maps from $V$ to $W$. Every linear map has $T(0) = 0$.
\end{definition}

\begin{theorem}[Linear map lemma]
If $v_1,\ldots,v_n$ is a basis of $V$ and $w_1,\ldots,w_n\in W$ are arbitrary, there is a unique linear map $T:V\to W$ with $Tv_k = w_k$.
\end{theorem}

With $(S+T)v = Sv+Tv$ and $(\lambda T)v = \lambda(Tv)$, $\mathcal{L}(V,W)$ is itself a vector space. For $T\in\mathcal{L}(U,V)$ and $S\in\mathcal{L}(V,W)$ the product $ST\in\mathcal{L}(U,W)$ is $(ST)u = S(Tu)$, defined only when $T$ maps into the domain of $S$. Products are associative, $(T_1T_2)T_3 = T_1(T_2T_3)$, have identities, $IT = TI = T$, and distribute, $(S_1+S_2)T = S_1T+S_2T$ and $T(S_1+S_2) = TS_1+TS_2$; but they are \emph{not} commutative.

\begin{theorem}[Rank--nullity]
For $T:V\to W$ with $V$ finite-dimensional, $\dim(\operatorname{null}T) + \dim(\operatorname{range}T) = \dim V$.
\end{theorem}

\begin{definition}[Matrix of a linear map]
An $m\times n$ matrix $A$ is a rectangular array of elements of $\mathbb{F}$ with entry $A_{j,k}$ in row $j$, column $k$. If $T\in\mathcal{L}(V,W)$, $v_1,\ldots,v_n$ is a basis of $V$ and $w_1,\ldots,w_m$ is a basis of $W$, the matrix of $T$ is the $m\times n$ matrix defined by $Tv_k = \sum_j A_{j,k}w_j$: column $k$ holds the components of $Tv_k$. Composition of maps corresponds to $(AB)_{\alpha\beta} = \sum_iA_{\alpha i}B_{i\beta}$.
\end{definition}

\section{Tensors}

\begin{definition}[Rank-$n$ tensor]
Let $\{\uvect{e}_i\}$ and $\{\uvect{e}'_i\}$ be two orthonormal bases of $\mathbb{R}^3$, related by the unique orthogonal matrix $R$ ($RR^T = R^TR = \mathbb{I}$) with $\uvect{e}'_i = \sum_jR_{ij}\uvect{e}_j$. A rank-$n$ tensor assigns $3^n$ components $T_{i_1\cdots i_n}$ to each orthonormal basis, related by
\[
    T'_{i_1\cdots i_n} = \sum_{j_1\cdots j_n}R_{i_1j_1}\cdots R_{i_nj_n}T_{j_1\cdots j_n}.
\]
For rank 2, $M'_{\alpha\beta} = \sum_{cd}R_{\alpha c}R_{\beta d}M_{cd}$, i.e.\ $M' = RMR^T$.
\end{definition}
